\clearpage
\chapter*{Erratum to Exposé no. 221}
\addcontentsline{toc}{chapter}{Erratum to Exposé no. 221}
\setcounter{footnote}{3}
\src{302}

\noindent\textbf{6. TDTE IV: Hilbert schemes.}

\begin{center}
[Bourbaki Seminar, vol. 13, 1960/61, no. 221, 28 pp.]
\end{center}

\noindent\textbf{Page 6, Theorem 2.1, third line.} --- For “it is necessary,”
read “it is necessary and sufficient.”

\noindent\textbf{Page 15, line 13.} --- For “$P\neq Q$,” read “$P<Q$.”

\noindent\textbf{Line 15.} --- For “If, on the contrary, $P=Q$,” read “In
the complementary case, $P(n)\geq Q(n)$ for large $n$.”

\noindent\textbf{Line 18, formula (*).} --- For “$<$” and “$>$,” read
“$\leq$” and “$\geq$.”

\noindent\textbf{Fifth line from the bottom.} --- For “$r$,” read “$Z$.”

\noindent\textbf{Page 18, Remark 3.9.} --- Let us point out that Hartshorne
has studied the connected components of Hilbert schemes over an algebraically
closed field. He proves that the $\operatorname{Hilb}^{P}_{r}$ are connected
and determines the pairs $(r,P)$ for which
$\operatorname{Hilb}^{P}_{r}\neq\varnothing$\footnote{Hartshorne (R.). ---
Connectedness of the Hilbert scheme (Thesis, Harvard, 1962).}.

\noindent\textbf{Page 23, fifth line from the bottom.} --- For “an
irreducible component,” read “of an irreducible component.”

\noindent\textbf{Pages 23 and 24, Remark 5.5.} --- In connection with Zappa's
example, let us point out that Mumford has just constructed an irreducible
component of the Hilbert scheme for $\mathbf P^3_{\mathbf C}$ (whose general
points represent nonsingular curves of degree 14 and genus 24) that is
nonreduced at its generic point. Blowing up one of the resulting curves, he
also obtains a regular projective scheme of dimension 3 over $\mathbf C$ whose
formal moduli scheme is nonreduced at its generic point or, equivalently,
whose local moduli variety in the sense of analytic geometry over $\mathbf C$
(cf. Cartan Seminar, vol. 13, 1960/61) is nonreduced at all its points.
